\addcontentsline{toc}{section}{Bibliography}
\section*{Bibliography}


 
[1] Sawtelle, S.D., 2018. Electromigration and Electrothermal Properties in Gold Nanostructures (Doctoral dissertation, Yale University).

[2] Moore, G.E., 1998. Cramming more components onto integrated circuits. Proceedings of the IEEE, 86(1), pp.82-85.


[3] Tang, C.L., 2005. Fundamentals of quantum mechanics: for solid state electronics and optics. Cambridge University Press.

[4] Nazarov, Y.V. and Blanter, Y.M., 2009. Quantum transport: introduction to nanoscience. Cambridge University Press.

[5] Combescot, R., 2022. Superconductivity: An Introduction. Cambridge University Press.

[6] Golubov, A.A., Wilhelm, F.K. and Zaikin, A.D., 1997. Coherent charge transport in metallic proximity structures. Physical Review B, 55(2), p.1123.

[7] de Bruijckere, J., 2020. Correlated Spin Phenomena in Molecular Quantum Transport Devices (Doctoral dissertation, Delft University of Technology).

[8] Gehring, P., Thijssen, J.M. and van der Zant, H.S., 2019. Single-molecule quantum-transport phenomena in break junctions. Nature Reviews Physics, 1(6), pp.381-396.

[9] Wang, X., Ouyang, Y., Jiao, L., Wang, H., Xie, L., Wu, J., Guo, J. and Dai, H., 2011. Graphene nanoribbons with smooth edges behave as quantum wires. Nature Nanotechnology, 6(9), pp.563-567.

[10] Jiao, L., Wang, X., Diankov, G., Wang, H. and Dai, H., 2010. Facile synthesis of high-quality graphene nanoribbons. Nature Nanotechnology, 5(5), pp.321-325.

[11] Ho, P.S. and Kwok, T., 1989. Electromigration in metals. Reports on Progress in Physics, 52(3), p.301.

[12] Baumans, X., 2018. Nano-engineering superconducting junctions by atomic migration.

[13] Likharev, K.K., 1979. Superconducting weak links. Reviews of Modern Physics, 51(1), p.101.

[14] Heida, J. P. (1998). Josephson currents in two dimensional mesoscopic ballistic conductors. [Thesis fully
internal (DIV), University of Groningen]. s.n.

[15] Martinis, J.M., Hilton, G.C., Irwin, K.D. and Wollman, D.A., 2000. Calculation of Tc in a normal-superconductor bilayer using the microscopic-based Usadel theory. Nuclear Instruments and Methods in Physics Research Section A: Accelerators, Spectrometers, Detectors and Associated Equipment, 444(1-2), pp.23-27.